% AI-produced by GPT-6 Astra (OpenAI), ChatGPT Pro mode, 18 November 2026.
% Working identifiers; the editor assigns the continuing Part B numbers.
\subsection{Faithfulness, good reduction, and the marked groupoid}

We fix the following convention.  A marked groupoid has a nonempty finite
set $I\subset X(K)$ of rational points as objects; its arrows are the
isomorphisms between the corresponding geometric fibre functors on finite
\'etale covers of $X_{\overline K}$.  The natural $G_K$-action fixes its
objects.  It acts on the actual arrows, not on their classes modulo inner
automorphisms.  Changing the geometric fibre functors conjugates this
action by a groupoid isomorphism and does not change its kernel.  The
singleton convention is already sufficient to test the expectation in
cote~142, scan~096, manuscript page~95.

\begin{claim}{B.N4.1}{Known}{The constant-field restriction on an outer kernel}
Let $k$ be a perfect field of characteristic zero and $X/k$ a hyperbolic
curve.  For the full geometric profinite group, the kernel of
$G_k\to\operatorname{Out}(\pi_1(X_{\overline k}))$ is contained in the
kernel of restriction $G_k\to G_{\mathbb Q}$, up to the usual choice of
an embedding of $\overline{\mathbb Q}$ in $\overline k$.  If $k$ is a
number field, the outer action, and hence every nonempty marked-groupoid
action, is faithful.
\end{claim}
\begin{proof}
This is Hoshi--Mochizuki, Theorem~C(ii), equivalently Corollary~6.2(ii),
specialized to configuration level $n=1$ and to the set of all primes.
Corollary~6.3(i) states the resulting injectivity for number fields and
$p$-adic local fields.  The source assumes a hyperbolic curve over a
perfect field and that every selected prime is invertible; these
hypotheses hold here.  Restriction $G_k\to G_{\mathbb Q}$ is injective
when $k/\mathbb Q$ is algebraic, which proves the number-field conclusion.
The kernel of an actual groupoid action lies in its outer-action kernel.
The proof of Corollary~6.2(ii) uses Belyi faithfulness on the
thrice-punctured line after Corollary~6.2(i).  It does not assert that
restriction $G_k\to G_{\mathbb Q}$ is injective when $k$ has positive
transcendence degree.  The attached author's PDF, printed pages~7 and
63--65, and the author's online copy were both opened.
\end{proof}

\begin{claim}{B.N4.2}{Proved}{Good reduction kills the actual marked action}
Let $R$ be a henselian DVR of residue characteristic zero, with fraction
field $F$.  Let $\overline{\mathcal X}/R$ be a smooth proper curve and
$\mathcal D\subset\overline{\mathcal X}$ a relative finite \'etale
divisor.  Put $\mathcal X=\overline{\mathcal X}\setminus\mathcal D$.
Suppose every point of a finite set $I\subset\mathcal X_F(F)$ extends
to an $R$-section of $\mathcal X$.  Then the inertia subgroup $I_R$ of
$G_F$ acts trivially on the full profinite groupoid
$\widehat\Pi_1(\mathcal X_{\overline F},I)$.
\end{claim}
\begin{proof}
Pass to the strict henselization $R^{\rm sh}$.  Its fraction field is
the maximal unramified extension of $F$ and its absolute Galois group
is $I_R$.  SGA~1, Expos\'e~XIII, Corollary~2.9 and the specialization
construction~2.10 apply to this smooth proper family with relative
normal-crossings boundary.  All primes are invertible in $R$, so the
prime-to-residue-characteristic restriction disappears.  They give the
equivalence of finite-cover categories between the geometric generic
fibre and the strictly local family (and the special fibre).  Equivalently,
every finite geometric cover is the restriction of a finite \'etale
cover of $\mathcal X_{R^{\rm sh}}$.

For an extending section $i$, the pullback of any such cover to
$\operatorname{Spec}R^{\rm sh}$ is a constant finite \'etale scheme.
Consequently its geometric generic fibre has trivial $I_R$-action.
The category of these extending covers is cofinal in all finite
geometric covers, and the same assertion holds simultaneously for each
of the finitely many fibre functors.  Naturality shows that inertia acts
identically on every isomorphism between those fibre functors.  Taking
the inverse limit proves triviality on every arrow of the profinite
groupoid.  In the singleton case this can also be seen from the based
specialization isomorphism: the rational section supplies its basepoint
transport, and $I_R$ acts trivially on the special-fibre group.
Thus an outer triviality statement has not been substituted for the
required actual-action statement.
\end{proof}

\begin{claim}{B.N4.3}{Proved}{Every finite marking over positive transcendence degree has a kernel}
Let $K$ be finitely generated over $\mathbb Q$, with
$\operatorname{trdeg}_{\mathbb Q}K>0$, and let $X/K$ be a smooth
hyperbolic curve.  For every nonempty finite set $I\subset X(K)$, the
actual action of $G_K$ on $\widehat\Pi_1(X_{\overline K},I)$ is not
faithful.  In fact its kernel contains a subgroup isomorphic to
$\widehat{\mathbb Z}(1)$.
\end{claim}
\begin{proof}
Choose a smooth integral model $S$ of $K$ over a number field.  Spreading
out the proper smooth completion, its finite reduced boundary, and the
finitely many marked points, and then shrinking $S$, gives a smooth
proper family with relative finite \'etale boundary and all marked
sections inside its complement.  There is a codimension-one point
$\eta$ of this nonempty positive-dimensional smooth scheme; its local
ring $R=\mathcal O_{S,\eta}$ is a DVR with residue characteristic zero
and fraction field $K$.  Its inertia subgroup embeds in $G_K$ after a
prolongation is fixed.  Tame Kummer theory gives
$I_R\simeq\widehat{\mathbb Z}(1)$, since the strict henselization has
value group $\mathbb Z$, separably closed residue field, and residue
characteristic zero.  In particular the subgroup is nontrivial: if
$\varpi$ is a uniformizer, adjoining $\sqrt[n]{\varpi}$ has ramification
index $n$ for every $n>1$.  Claim~B.N4.2 puts this whole subgroup in the
marked-action kernel.  Every comparison uses restriction from the same
chosen algebraic closure of $K$; changing the prolongation merely
conjugates the displayed subgroup.
\end{proof}

\begin{claim}{B.N4.4}{False as stated}{The moduli condition does not imply uniqueness}
The uniqueness expectation of cote~142, scan~096 is false for finite
markings, even when the field of definition is finite over the field of
moduli.  A counterexample is
\[
 K=\mathbb Q(t),\qquad
 X=\mathbb P^1_K\setminus\{0,1,t,\infty\},\qquad I=\{2\}.
\]
The inertia at $t=3$ is a nontrivial subgroup of the kernel of the actual
marked action.  The assertion becomes true for number fields, and for
finite markings positive transcendence degree always prevents it, as
specified in Claims~B.N4.1 and B.N4.3.
\end{claim}
\begin{proof}
The curve has type $(0,4)$ and is hyperbolic.  Set
\[
 j(t)=256\frac{(1-t+t^2)^3}{t^2(1-t)^2}.
\]
The six cross-ratio transforms of the unordered four-point set are
$t,1-t,1/t,1/(1-t),t/(t-1),(t-1)/t$.  They form its generic orbit,
and $j$ is their common invariant.  The rational map $t\mapsto j(t)$
has degree six (its coprime numerator and denominator have degrees six
and four).  Thus the fixed field of this order-six group is exactly
$\mathbb Q(j)$, the field of moduli of the unlabelled punctured curve,
and $[K:\mathbb Q(j)]=6$.  If the punctures are ordered, the moduli
field is already $K$; adding the rational marking does not invalidate
the finiteness assertion.

Over $R=\mathbb Q[t]_{(t-3)}$, the sections $0,1,t,\infty$ are disjoint
and the section $2$ avoids them.  Claim~B.N4.2 therefore kills the
inertia at $t-3$ in the actual based action.  This inertia is nontrivial
because $K(\sqrt{t-3})/K$ is ramified of degree two at that valuation;
indeed all Kummer degrees occur after the unramified constant extension.
This proves a counterexample in the full profinite group and does not
rely on a finite nilpotent approximation.  The program
\texttt{n4\_moduli\_certificate.py} checks the six rational invariants,
the polynomial degree and coprimality, and the specialization of the
marked divisor.  The ramification and trivial-action assertions are the
preceding proofs, not claims made by that finite calculation.
\end{proof}

\begin{claim}{B.N4.5}{Proved}{The quotient that good-reduction inertia actually supplies}
Let a marked curve as above spread to a smooth relative pair over a
smooth integral characteristic-zero scheme $S$, with function field
$K$.  Let $N_S$ be the closed normal subgroup of $G_K$ generated by
the inertia groups at all codimension-one points of $S$.  Then
\[
 G_K/N_S\simeq\pi_1(S),\qquad
 \ker(\rho_I)/N_S=\ker(\rho_{S,I}),
\]
where $\rho_{S,I}:\pi_1(S)\to\operatorname{Aut}
(\widehat\Pi_1(X_{\overline K},I))$ is the actual family monodromy.
Thus uniqueness modulo $N_S$ is precisely injectivity of this family
monodromy; it is not a consequence of the moduli-field condition alone.
\end{claim}
\begin{proof}
The restriction functor from finite \'etale covers of normal $S$ to
finite separable extensions of $K$ is fully faithful, and its image
consists of the extensions whose normalizations are unramified over
$S$.  Since $S$ is regular, purity of the branch locus identifies this
image with extensions unramified at every codimension-one point.
Equivalently the corresponding finite quotients of $G_K$ are exactly
those killing all of the indicated inertia subgroups.  This proves the
first displayed isomorphism by the universal property of the profinite
quotient.  The smooth pair supplies finite locally constant systems of
geometric covers and of their fibres at the marked sections; hence its
groupoid monodromy is defined on $\pi_1(S)$.  Restriction to the generic
point is $\rho_I$.  Kernels in this factorization give the second
identity.  Purity is used here in its usual regular excellent form:
Stacks Project, Lemma~58.21.4 (Tag~0BMB), applied to the finite
normalization.  This result gives a factorization and the exact
remaining monodromy kernel, not an unproved assertion that the latter
is always zero.
\end{proof}

\begin{claim}{B.N4.6}{Proved}{A faithful corrected monodromy in the counterexample}
For the curve and marking of Claim~B.N4.4, put
$S=\mathbb P^1_{\mathbb Q}\setminus\{0,1,2,\infty\}$.
Its actual family monodromy
\[
 \pi_1(S)\longrightarrow
 \operatorname{Aut}\bigl(\pi_1(X_{\overline{\mathbb Q(t)}},2)\bigr)
\]
is injective.  Consequently the kernel of the $G_{\mathbb Q(t)}$-action
is exactly $N_S$ in Claim~B.N4.5.
\end{claim}
\begin{proof}
Write $C=\mathbb P^1\setminus\{0,1,\infty\}$ and
$Y=(C\times C)\setminus\operatorname{diag}$.  Work first over
$\overline{\mathbb Q}$ with the point $(x,t)=(2,3)$.  In
$E=\pi_1(Y)$ let $F$ be the fibre subgroup of the second projection
and $H$ the fibre subgroup of the first projection.  Both are free
profinite groups of rank three and both inject in $E$.  Here full
profinite fibre exactness can be verified directly.  Over the complex
numbers the projection is a locally trivial punctured-surface fibration
with base having free discrete group of rank two.  Its exact sequence
splits because that quotient is free.  For every finite quotient of
the finitely generated free fibre group, choose a characteristic
finite-index subgroup inside its kernel.  The resulting finite fibre
group $Q$ has finite automorphism group; the semidirect product maps
to $Q\rtimes\operatorname{Aut}(Q)$, proving that every finite fibre
quotient is detected in a finite quotient of the total group.  Right
exactness of completion proves the assertion.  Riemann existence and
invariance under algebraically closed characteristic-zero extension
identify these completions with the groups in use.

The geometric monodromy of the marked family is the conjugation of
$H$ on $F$: the map $S\to Y$, $t\mapsto(2,t)$, is exactly the fibre
inclusion for the first projection and supplies the section for the
second projection after restricting its base to $S$.  The intersection
$F\cap H$ is the kernel of
$H=\pi_1(C\setminus\{2\})\to\pi_1(C)$.  It contains the cusp
inertia group $J$ at $2$ and every $H$-conjugate of $J$.

The cusp facts needed next are Stix (2012), Proposition~16(1)--(3):
distinct geometric cusp inertia groups have trivial intersection,
their geometric normalizers equal themselves, and an abelian subgroup
meeting such inertia nontrivially lies in it.  They hold for this
hyperbolic curve.  A topological generator $c$ of $J$ has
$C_H(c)=J$: its centralizer normalizes the procyclic closure of
$\langle c\rangle$, so lies in $J$, and the reverse inclusion is clear.
There exists $h\in H\setminus J$, because $H$ has rank three and is
not procyclic.  Then $J\cap hJh^{-1}=1$.  Any element of $H$ acting
trivially on $F$ centralizes both $c$ and $hch^{-1}$, and hence lies in
this trivial intersection.  Geometric monodromy is injective.

Finally filling the moving puncture gives an equivariant surjection
$F\twoheadrightarrow\pi_1(C_{\overline{\mathbb Q}},2)$.
The induced arithmetic monodromy on this quotient is the constant
curve's $G_{\mathbb Q}$-action.  Claim~B.N4.1 makes its outer action
faithful.  An element of the full $\pi_1(S)$-kernel therefore has
trivial image in $G_{\mathbb Q}$, so belongs to the geometric subgroup,
where injectivity was just proved.  This establishes the stated
arithmetic injection and then the kernel equality.
\end{proof}

\begin{claim}{B.N4.7}{Proved}{The constant example and the scope of image conditions}
For a constant hyperbolic curve $X_0/\mathbb Q$ with constant nonempty
finite marking $I_0$, its base change to $K=\mathbb Q(t)$ has marked
action factoring through $G_K\to G_{\mathbb Q}$; its kernel is exactly
the kernel of that restriction.  This explains the isotrivial
obstruction, but Claims~B.N4.3--B.N4.4 show that eliminating
isotriviality does not eliminate good-reduction inertia.

The additional relations of Nakamura--Schneps (2000), Theorem~1.1 and
Definition~1.1, define $\mathrm I\Gamma\subset\widehat{GT}$;
Theorem~1.2 makes it a subgroup containing the Galois image.  None of
these statements identifies it with all of $\widehat{GT}$.
Likewise Hoshi--Mochizuki--Tsujimura (2025), Definition~3.3(ii),(iii),(v),
imposes the COF and RGC properties on closed subgroups before denoting
them by $\mathrm{BGT}$.  These are hypotheses on the selected subgroups,
not assertions about every element of the ambient group.
\end{claim}
\begin{proof}
In the constant example, covers and the constant fibre functors descend
from the algebraic closure of the constant field.  Thus the action
factors as asserted.  Faithfulness over $\mathbb Q$ from Claim~B.N4.1
identifies its kernel.  The source distinctions in the second paragraph
follow by reading the stated definitions and theorems in the opened
primary texts.  They cannot be used to prove that the necessary
conditions a), b) of cote~142, scan~095 characterize the whole ambient
automorphism group.  Those conditions concern existence of a Galois
element; the counterexample in Claim~B.N4.4 concerns uniqueness of an
element already known to act.  The two assertions are logically
separate.  The Belyi argument in Claim~B.N4.1 proves faithfulness for
the constant-field Galois quotient and does not remove geometric
inertia over a field of positive transcendence degree.
\end{proof}

\paragraph{Check record.}
Claims B.N4.2--B.N4.7 were checked only in this writing by GPT-6 Astra
(OpenAI, ChatGPT Pro mode), 18 November 2026: direct re-proof of the
specialization and fibre-functor comparisons, independent construction
of the counterexample, and the explicit symbolic certificate specified
above.  The separately selected internal reasoning-effort value is not
exposed.  B.N4.1 is an opened published theorem with its hypotheses
matched here.  No human or different-model-family verification is
claimed.  The field-of-moduli condition and the finite-marking
convention are stated explicitly; no assertion about arbitrary infinite
marking sets is silently inferred.
